\chapter{Observers and Kalman Filtering}

\section{State Estimation and Luenberger Observers}
In the previous chapter, we designed state feedback controllers under the assumption that the full state vector $x(t)$ is directly measurable. In practical applications, however, this assumption rarely holds. Measuring every state variable can be cost-prohibitive, technically challenging due to sensor physical limitations, or corrupted by high-frequency noise. Instead, we typically only have access to a vector of control inputs $u(t)$ and a vector of measured outputs $y(t) \in \R^p$, where $p < n$.

To overcome this, we design a dynamic system called a \textbf{state observer} (or estimator) that runs in real-time in the control computer. The observer processes the inputs $u(t)$ and outputs $y(t)$ to generate an estimate of the state vector, denoted by $\hat{x}(t) \in \R^n$.

Consider the continuous-time LTI system:
\begin{align}
    \dot{x}(t) &= A x(t) + B u(t) \label{eq:obs_plant_state} \\
    y(t) &= C x(t) + D u(t) . \label{eq:obs_plant_output}
\end{align}

A naive attempt to estimate $x(t)$ would be to simulate the plant model directly: $\dot{\hat{x}} = A \hat{x} + B u$. However, any initial state mismatch ($x(0) \neq \hat{x}(0)$) or model inaccuracy would cause the estimate to drift over time. 
The step-by-step structural evolution of the observer is illustrated in Figures \ref{fig:naive_estimator}, \ref{fig:correction_feedback}, and \ref{fig:closed_loop_observer} (collectively shown in Figure \ref{fig:observer_evolution}).

\begin{figure}[H]
    \centering
    \begin{minipage}{0.31\textwidth}
        \centering
        \includegraphics[width=\textwidth]{state_observer2.png}
        \caption{Naive open-loop estimator structure.}
        \label{fig:naive_estimator}
    \end{minipage}
    \hfill
    \begin{minipage}{0.31\textwidth}
        \centering
        \includegraphics[width=\textwidth]{state_observer3.png}
        \caption{Introducing correction feedback block.}
        \label{fig:correction_feedback}
    \end{minipage}
    \hfill
    \begin{minipage}{0.31\textwidth}
        \centering
        \includegraphics[width=\textwidth]{state_observer4.png}
        \caption{Final closed-loop observer dynamics.}
        \label{fig:closed_loop_observer}
    \end{minipage}
    \caption{Step-by-step structural evolution of a state observer from open-loop simulation to closed-loop estimation.}
    \label{fig:observer_evolution}
\end{figure}

To prevent this drift, David Luenberger proposed feeding back the measurement error (the difference between the actual output $y(t)$ and the estimated output $\hat{y}(t) = C \hat{x}(t) + D u(t)$) to correct the state estimates:
\begin{equation}
    \dot{\hat{x}}(t) = A \hat{x}(t) + B u(t) + L \left( y(t) - C \hat{x}(t) - D u(t) \right) , \label{eq:luenberger_obs}
\end{equation}
where $L \in \R^{n \times p}$ is the observer gain matrix. The completed Luenberger observer structure is shown in Figure \ref{fig:state_observer}.

\begin{figure}[H]
    \centering
    \includegraphics[width=9cm]{state_observer.png}
    \caption{Luenberger state observer block diagram.}
    \label{fig:state_observer}
\end{figure}

\subsection{Observer Error Dynamics}
Let us define the estimation error vector as:
\begin{equation}
    e(t) = x(t) - \hat{x}(t) . \label{eq:obs_error_def}
\end{equation}
To analyze how the error behaves over time, we differentiate \eqref{eq:obs_error_def}:
\begin{align}
    \dot{e}(t) &= \dot{x}(t) - \dot{\hat{x}}(t) \nonumber \\
    &= (A x + B u) - \left( A \hat{x} + B u + L (C x + D u - C \hat{x} - D u) \right) \nonumber \\
    &= A (x - \hat{x}) - L C (x - \hat{x}) \nonumber \\
    &= (A - L C) e(t) . \label{eq:obs_error_dynamics}
\end{align}
The error dynamics are governed by the autonomous system matrix $A - L C$. If $A - L C$ is Hurwitz, the estimation error will converge to zero asymptotically: $\lim_{t \to \infty} e(t) = 0$, regardless of the initial mismatch $e(0)$ and the control input $u(t)$.
The continuous-time observer implementation and its error trajectories are shown in Figures \ref{fig:observer_continuous} and \ref{fig:observer_example} (collectively shown in Figure \ref{fig:observer_cont_sim}).

\begin{figure}[H]
    \centering
    \begin{minipage}{0.48\textwidth}
        \centering
        \includegraphics[width=\textwidth]{state_observer_continuous.png}
        \caption{Continuous-time implementation of Luenberger observer.}
        \label{fig:observer_continuous}
    \end{minipage}
    \hfill
    \begin{minipage}{0.48\textwidth}
        \centering
        \includegraphics[width=\textwidth]{state_observer_example.png}
        \caption{Observer estimation error trajectories during initial transient.}
        \label{fig:observer_example}
    \end{minipage}
    \caption{Continuous-time state observer structure and numerical simulation trajectories.}
    \label{fig:observer_cont_sim}
\end{figure}


\begin{mytheorem}{Arbitrary Observer Pole Placement}
The eigenvalues of the observer matrix $A - L C$ can be placed at arbitrary, symmetric locations in the complex plane by selecting an appropriate observer gain matrix $L$ if and only if the system pair $(A, C)$ is observable.
\end{mytheorem}

Mathematically, this is the dual of the state feedback problem. By taking the transpose:
\begin{equation}
(A - L C)^\trp = A^\trp - C^\trp L^\trp ,
\end{equation}
we see that designing $L$ is equivalent to designing a state feedback gain $K = L^\trp$ for the dual pair $(A^\trp, C^\trp)$. Thus, Ackermann's formula can be applied directly to compute $L$:
\begin{equation}
L = \left( \begin{bmatrix} 0 & \dots & 0 & 1 \end{bmatrix} \mathcal{O}_{dual}^{-1} \alpha_{obs}(A^\trp) \right)^\trp ,
\end{equation}
where $\mathcal{O}_{dual}$ is the controllability matrix of the dual system (which is the transpose of the observability matrix $\mathcal{O}$ of the original system).

\section{The Separation Principle}
We now combine the state feedback controller and the Luenberger observer. Since we do not have direct access to $x(t)$, we feed back the estimated state $\hat{x}(t)$ instead:
\begin{equation}
    u(t) = -K \hat{x}(t) + r(t) . \label{eq:sf_obs_law}
\end{equation}
This combined controller-observer structure is called an \textbf{output feedback controller}. Since the controller relies on estimates and the estimator relies on the control input, they are coupled. We must prove that this coupled system remains stable and that the controller and observer gains can be designed independently.

\begin{mytheorem}{The Separation Principle}
For the LTI system \eqref{eq:obs_plant_state}--\eqref{eq:obs_plant_output} under the estimator-based control law \eqref{eq:sf_obs_law}, the closed-loop eigenvalues are exactly the union of the controller eigenvalues (eigenvalues of $A - B K$) and the observer eigenvalues (eigenvalues of $A - L C$).
\end{mytheorem}

\begin{proof}
Let us choose the state vector of the augmented system to be $x_{aug}(t) = \begin{bmatrix} x(t)^\trp & e(t)^\trp \end{bmatrix}^\trp$. We express the closed-loop dynamics in these coordinates.
First, substitute $\hat{x} = x - e$ into the control law \eqref{eq:sf_obs_law}:
\begin{equation}
u(t) = -K (x(t) - e(t)) + r(t) = -K x(t) + K e(t) + r(t) .
\end{equation}
Substitute this into the plant state equation:
\begin{align}
    \dot{x}(t) &= A x(t) + B \left( -K x(t) + K e(t) + r(t) \right) \nonumber \\
    &= (A - B K) x(t) + B K e(t) + B r(t) . \label{eq:sep_x}
\end{align}
The error dynamics remain unaffected by the control input, as shown in \eqref{eq:obs_error_dynamics}:
\begin{equation}
    \dot{e}(t) = (A - L C) e(t) . \label{eq:sep_e}
\end{equation}
Combining \eqref{eq:sep_x} and \eqref{eq:sep_e} into a single state-space system yields:
\begin{equation}
    \begin{bmatrix} \dot{x}(t) \\ \dot{e}(t) \end{bmatrix} = \begin{bmatrix} A - B K & B K \\ 0 & A - L C \end{bmatrix} \begin{bmatrix} x(t) \\ e(t) \end{bmatrix} + \begin{bmatrix} B \\ 0 \end{bmatrix} r(t) . \label{eq:sep_augmented}
\end{equation}
To find the closed-loop poles, we compute the characteristic polynomial of the augmented system matrix:
\begin{equation}
\det \left( s I - \begin{bmatrix} A - B K & B K \\ 0 & A - L C \end{bmatrix} \right) = 0 .
\end{equation}
Because the matrix is block upper-triangular, its determinant is the product of the determinants of the diagonal blocks:
\begin{equation}
\det \left( s I - (A - B K) \right) \det \left( s I - (A - L C) \right) = 0 .
\end{equation}
This proves that the closed-loop poles are precisely the poles of $A - B K$ together with the poles of $A - L C$. The design of the controller and the design of the observer are completely \textbf{separated}! We can design $K$ to meet transient specifications assuming full state feedback, and design $L$ to ensure the estimation error converges rapidly, without any risk of coupled instability.
\end{proof}

In practice, a common rule of thumb is to design the observer poles to be $2$ to $10$ times faster (more negative) than the controller poles. This ensures that the state estimates converge rapidly, providing accurate information to the controller during transients.

\section{The Kalman Filter}
The Luenberger observer is deterministic: it assumes perfect model parameters and noise-free measurements. In reality, control systems are subjected to stochastic process disturbances (e.g., wind gusts on an aircraft) and sensor measurement noise.

The closed-loop loop structure under LQR control and Kalman filtering is shown in Figure \ref{fig:state_observer_kf}.

\begin{figure}[H]
    \centering
    \includegraphics[width=9cm]{state_observer_kf.png}
    \caption{Kalman filter state estimation and feedback loop structure.}
    \label{fig:state_observer_kf}
\end{figure}

The \textbf{Kalman Filter}, introduced by Rudolf E. Kálmán in 1960, is the mathematically optimal estimator for linear systems subjected to Gaussian white noise.

Rudolf E. Kalman and his landmark 1960 optimal filtering paper are shown in Figures \ref{fig:Rudolf_Kalman} and \ref{fig:paper_1960_kf} (collectively shown in Figure \ref{fig:kalman_history}).

\begin{figure}[H]
    \centering
    \begin{minipage}{0.45\textwidth}
        \centering
        \includegraphics[width=3.5cm]{Rudolf_Kalman.jpg}
        \caption{Rudolf E. Kálmán (1930--2016).}
        \label{fig:Rudolf_Kalman}
    \end{minipage}
    \hfill
    \begin{minipage}{0.45\textwidth}
        \centering
        \includegraphics[width=5.5cm]{paper_1960_kf.png}
        \caption{Kalman's landmark 1960 paper published in ASME Trans.}
        \label{fig:paper_1960_kf}
    \end{minipage}
    \caption{Historical founder and publication of the Kalman Filter.}
    \label{fig:kalman_history}
\end{figure}


Consider the discrete-time stochastic LTI system:
\begin{align}
    x_{k+1} &= A x_k + B u_k + w_k \label{eq:kf_state} \\
    y_k &= C x_k + v_k . \label{eq:kf_output}
\end{align}

This stochastic system structure is shown in Figure \ref{fig:modelo_ss_discreto_estocastico}.

\begin{figure}[H]
    \centering
    \includegraphics[width=9cm]{modelo_ss_discreto_estocastico.png}
    \caption{Discrete-time stochastic state-space block diagram showing process and measurement noise injection.}
    \label{fig:modelo_ss_discreto_estocastico}
\end{figure}

where:
- $w_k \in \R^n$ is the \textbf{process noise}, representing random unmodeled disturbances.
- $v_k \in \R^p$ is the \textbf{measurement noise}, representing sensor inaccuracies.

We assume that $w_k$ and $v_k$ are zero-mean, mutually independent white Gaussian random processes with known covariance matrices:
\begin{equation}
\mathbb{E}[w_k] = 0, \quad \mathbb{E}[v_k] = 0 ,
\end{equation}
\begin{equation}
\mathbb{E}[w_k w_j^\trp] = W \delta_{kj}, \quad \mathbb{E}[v_k v_j^\trp] = V \delta_{kj}, \quad \mathbb{E}[w_k v_j^\trp] = 0 ,
\end{equation}
where $W = W^\trp \ge 0$, $V = V^\trp \succ 0$, and $\delta_{kj}$ is the Kronecker delta.

\subsection{Gaussian Distributions}
The process and measurement noise are modeled as Gaussian random variables. The scalar probability density functions under different means $\mu$ and standard deviations $\sigma$ are plotted in Figures \ref{fig:gauss_1}, \ref{fig:gauss_2}, and \ref{fig:gauss_3} (collectively shown in Figure \ref{fig:gaussiana_escalar}), while the multidimensional Gaussian distribution is illustrated in Figure \ref{fig:2d_gaussian}.

\begin{figure}[H]
    \centering
    \begin{minipage}{0.31\textwidth}
        \centering
        \includegraphics[width=\textwidth]{gaussiana_escalar.png}
        \caption{$\mu=0$, $\sigma=1$.}
        \label{fig:gauss_1}
    \end{minipage}
    \hfill
    \begin{minipage}{0.31\textwidth}
        \centering
        \includegraphics[width=\textwidth]{gaussiana_escalar2.png}
        \caption{$\mu=2$, $\sigma=1$.}
        \label{fig:gauss_2}
    \end{minipage}
    \hfill
    \begin{minipage}{0.31\textwidth}
        \centering
        \includegraphics[width=\textwidth]{gaussiana_escalar3.png}
        \caption{$\mu=0$, $\sigma=2$.}
        \label{fig:gauss_3}
    \end{minipage}
    \caption{Scalar Gaussian probability density functions under various parameters.}
    \label{fig:gaussiana_escalar}
\end{figure}

\begin{figure}[H]
    \centering
    \includegraphics[width=8cm]{2d_gaussian.png}
    \caption{Probability density surface of a two-dimensional multivariate Gaussian distribution.}
    \label{fig:2d_gaussian}
\end{figure}


\subsection{State and Covariance Definitions}
The Kalman filter maintains the conditional mean and covariance of the state. We define:
- $\hat{x}_{k|k-1}$: The \textbf{a priori} state estimate at step $k$, given all measurements up to step $k-1$.
- $\hat{x}_{k|k}$: The \textbf{a posteriori} state estimate at step $k$, given all measurements up to step $k$.
- $P_{k|k-1}$: The \textbf{a priori} estimation error covariance matrix:
  \begin{equation}
P_{k|k-1} = \mathbb{E}\left[ (x_k - \hat{x}_{k|k-1})(x_k - \hat{x}_{k|k-1})^\trp \right] .
  \end{equation}
- $P_{k|k}$: The \textbf{a posteriori} estimation error covariance matrix:
  \begin{equation}
P_{k|k} = \mathbb{E}\left[ (x_k - \hat{x}_{k|k})(x_k - \hat{x}_{k|k})^\trp \right] .
  \end{equation}

The filter operates in two alternating phases: \textbf{Prediction (Time Update)} and \textbf{Correction (Measurement Update)}.

\subsection{Prediction Phase (Time Update)}
In the prediction phase, the filter projects the state and covariance ahead in time using the system dynamics.
\begin{enumerate}
    \item \textbf{State Prediction}:
    Since the noise $w_k$ is zero-mean and uncorrelated with previous states, the optimal prediction of the state is:
    \begin{equation}
        \hat{x}_{k|k-1} = A \hat{x}_{k-1|k-1} + B u_{k-1} , \label{eq:kf_pred_state}
    \end{equation}
    \item \textbf{Covariance Prediction}:
    Let us define the a priori error $e_{k|k-1} = x_k - \hat{x}_{k|k-1}$. Substituting the dynamics \eqref{eq:kf_state} and prediction \eqref{eq:kf_pred_state}:
    \begin{align}
        e_{k|k-1} &= (A x_{k-1} + B u_{k-1} + w_{k-1}) - (A \hat{x}_{k-1|k-1} + B u_{k-1}) \nonumber \\
&= A e_{k-1|k-1} + w_{k-1} .
    \end{align}
    The a priori covariance is:
    \begin{align}
        P_{k|k-1} &= \mathbb{E}[e_{k|k-1} e_{k|k-1}^\trp] \nonumber \\
        &= \mathbb{E}\left[ (A e_{k-1|k-1} + w_{k-1})(A e_{k-1|k-1} + w_{k-1})^\trp \right] \nonumber \\
&= A \mathbb{E}[e_{k-1|k-1} e_{k-1|k-1}^\trp] A^\trp + \mathbb{E}[w_{k-1} w_{k-1}^\trp] + A \mathbb{E}[e_{k-1|k-1} w_{k-1}^\trp] + \mathbb{E}[w_{k-1} e_{k-1|k-1}^\trp] A^\trp .
    \end{align}
    Since the process noise $w_{k-1}$ is white and uncorrelated with the past error $e_{k-1|k-1}$, the cross terms vanish. This simplifies to:
    \begin{equation}
        P_{k|k-1} = A P_{k-1|k-1} A^\trp + W . \label{eq:kf_pred_cov}
    \end{equation}
\end{enumerate}

\subsection{Correction Phase (Measurement Update)}
When a new measurement $y_k$ becomes available, we correct our prediction using a Luenberger-like update:
\begin{equation}
    \hat{x}_{k|k} = \hat{x}_{k|k-1} + K_k \left( y_k - C \hat{x}_{k|k-1} \right) , \label{eq:kf_corr_state}
\end{equation}
where $K_k \in \R^{n \times p}$ is the \textbf{Kalman Gain} to be determined.

To find the optimal gain $K_k$, we analyze the a posteriori error $e_{k|k} = x_k - \hat{x}_{k|k}$:
\begin{align}
    e_{k|k} &= x_k - \left( \hat{x}_{k|k-1} + K_k (C x_k + v_k - C \hat{x}_{k|k-1}) \right) \nonumber \\
&= (I - K_k C) e_{k|k-1} - K_k v_k .
\end{align}
The a posteriori covariance is:
\begin{align}
    P_{k|k} &= \mathbb{E}[e_{k|k} e_{k|k}^\trp] \nonumber \\
    &= \mathbb{E}\left[ \left( (I - K_k C) e_{k|k-1} - K_k v_k \right) \left( (I - K_k C) e_{k|k-1} - K_k v_k \right)^\trp \right] \nonumber \\
    &= (I - K_k C) \mathbb{E}[e_{k|k-1} e_{k|k-1}^\trp] (I - K_k C)^\trp + K_k \mathbb{E}[v_k v_k^\trp] K_k^\trp \nonumber \\
&\quad - (I - K_k C) \mathbb{E}[e_{k|k-1} v_k^\trp] K_k^\trp - K_k \mathbb{E}[v_k e_{k|k-1}^\trp] (I - K_k C)^\trp .
\end{align}
Since the sensor noise $v_k$ is uncorrelated with the a priori error $e_{k|k-1}$, the cross terms vanish:
\begin{equation}
    P_{k|k} = (I - K_k C) P_{k|k-1} (I - K_k C)^\trp + K_k V K_k^\trp . \label{eq:kf_joseph_form}
\end{equation}
Equation \eqref{eq:kf_joseph_form} is known as the \textbf{Joseph Form} of the covariance update, which is numerically stable and holds for any choice of gain $K_k$.

\subsection{Derivation of the Optimal Kalman Gain}
The optimal Kalman gain $K_k$ is defined as the gain that minimizes the total mean-squared estimation error, which is the sum of the variances of the state estimation errors. This is mathematically represented by the trace of the covariance matrix $P_{k|k}$:
\begin{equation}
\min_{K_k} \tr(P_{k|k}) .
\end{equation}
Expanding \eqref{eq:kf_joseph_form}:
\begin{equation}
P_{k|k} = P_{k|k-1} - K_k C P_{k|k-1} - P_{k|k-1} C^\trp K_k^\trp + K_k \left( C P_{k|k-1} C^\trp + V \right) K_k^\trp .
\end{equation}
Taking the trace of both sides and using the cyclic properties of the trace (e.g., $\tr(A) = \tr(A^\trp)$):
\begin{equation}
\tr(P_{k|k}) = \tr(P_{k|k-1}) - 2 \tr(K_k C P_{k|k-1}) + \tr\left( K_k (C P_{k|k-1} C^\trp + V) K_k^\trp \right) .
\end{equation}
To find the minimum, we take the matrix derivative with respect to $K_k$ and set it to zero:
\begin{equation}
\frac{\partial \tr(P_{k|k})}{\partial K_k} = -2 P_{k|k-1} C^\trp + 2 K_k \left( C P_{k|k-1} C^\trp + V \right) = 0 .
\end{equation}
Solving for $K_k$ yields the \textbf{optimal Kalman Gain}:
\begin{equation}
    K_k = P_{k|k-1} C^\trp \left( C P_{k|k-1} C^\trp + V \right)^{-1} . \label{eq:kf_optimal_gain}
\end{equation}

Substituting this optimal gain back into the Joseph form \eqref{eq:kf_joseph_form} simplifies the expression for the a posteriori covariance:
\begin{align}
    P_{k|k} &= (I - K_k C) P_{k|k-1} (I - K_k C)^\trp + K_k V K_k^\trp \nonumber \\
&= P_{k|k-1} - K_k C P_{k|k-1} - P_{k|k-1} C^\trp K_k^\trp + K_k (C P_{k|k-1} C^\trp + V) K_k^\trp .
\end{align}
Using $K_k (C P_{k|k-1} C^\trp + V) = P_{k|k-1} C^\trp$, the last two terms cancel:
\begin{equation}
    P_{k|k} = (I - K_k C) P_{k|k-1} . \label{eq:kf_corr_cov}
\end{equation}

\subsection{Summary of the Kalman Filter Algorithm}
The complete recursive equations of the Kalman filter are summarized in Table~\ref{tab:kalman_filter}.

\begin{table}[H]
\centering
\caption{The Discrete-Time Kalman Filter Algorithm}
\label{tab:kalman_filter}
\begin{tabular}{lll}
\toprule
\textbf{Phase} & \textbf{Equation} & \textbf{Description} \\
\midrule
\textbf{Prediction} & $\hat{x}_{k|k-1} = A \hat{x}_{k-1|k-1} + B u_{k-1}$ & Project the state ahead \\
                    & $P_{k|k-1} = A P_{k-1|k-1} A^\trp + W$ & Project the error covariance ahead \\
\midrule
\textbf{Correction} & $K_k = P_{k|k-1} C^\trp (C P_{k|k-1} C^\trp + V)^{-1}$ & Compute the optimal Kalman Gain \\
                    & $\hat{x}_{k|k} = \hat{x}_{k|k-1} + K_k (y_k - C \hat{x}_{k|k-1})$ & Update estimate with measurement \\
                    & $P_{k|k} = (I - K_k C) P_{k|k-1}$ & Update the error covariance \\
\bottomrule
\end{tabular}
\end{table}

Under the assumptions of stabilizability of $(A, B)$ and detectability of $(A, C)$, the error covariance converges to a steady-state value $P_{k|k-1} \to P$ as $k \to \infty$. This steady-state covariance satisfies the discrete Riccati equation:
\begin{equation}
P = A P A^\trp - A P C^\trp (C P C^\trp + V)^{-1} C P A^\top + W ,
\end{equation}
yielding a constant optimal filter gain $K = P C^\trp (C P C^\trp + V)^{-1}$. This steady-state estimator is identical in structure to a discrete Luenberger observer, but its gain is optimally tuned to balance process disturbances and sensor noise.

\section{Homework: State Observer for a Physical System}
\begin{figure}[H]
    \centering
    \includegraphics[width=5cm]{pendulum.png}
    \caption{State estimation task on an inverted pendulum system.}
    \label{fig:homework_observer_pendulum}
  \end{figure}
  
  \begin{exercise}
  Select a physical system of your interest (e.g., the inverted pendulum shown in Figure \ref{fig:homework_observer_pendulum}, a quadcopter, or a bioreactor). Design a Luenberger observer in continuous-time and verify that the observer poles are faster than the closed-loop controller poles. Under initial state mismatches, plot the estimated states alongside the actual states to demonstrate convergence of the estimation error.
  \end{exercise}
